\begin{figure}[h!t]

    \centering
    
    \begin{tikzpicture}
        
        % layout
        \def\ARC{1.2cm}     % radius of the arc used to indicate angles
        
        % parameters
        \def\Px{1.75}       % x coordinate of point P
        \def\Py{0.00}       % y coordinate of point P
%
%================================================
%                                
        % Cartesian coordinate system
        \begin{scope}[xshift=0cm,yshift=0cm]
\ifdebug
            % grid lines
            \path[draw,grid=0.25cm] (-3.50,-3.00) grid (3.00,2.75);
\fi
            
            \begin{scope}[rotate=47.5]
                % objects
                \path coordinate (O) at (0,0);          % origin O
                \path coordinate (P) at (\Px,\Py);      % point P
               % reference axes
                \draw[main_axis] (-3.60,0) -- (3.25,0) node[below right] {$$};
                % origin
                \path node at (O) {$\mathcal{O}$};
                % tick marks at unit intervals
                \foreach \x in {-3,...,3}
                {
                    \path[tick] (\x,-2pt) -- (\x,+2pt);
                }
                \foreach \x in {-3,...,-1}
                {
                    \path node[above left, inner sep=2pt]
                        at (\x,0) {$\x$};
                }
                \foreach \x in {1,...,3}
                {
                    \path node[above left, inner sep=2pt]
                        at (\x,0) {$\x$};
                }
                    
                % indicate positive and negative half-line
                \path node at ( 3.45,0.00) {$\mathbf{+}$};
                \path node at (-3.75,0.00) {$\mathbf{-}$};                
                % point P
                \filldraw[blue] (P) circle (2pt);
                \path node[text=blue,below right, inner sep=1pt]
                    at (P) {$\mathcal{P}=(d_\mathcal{P})$};
            
            \end{scope}
            
            % label sub figure
            \path node[above right] at (-3.50,-3.00) {[\textbf{a}]};
                        
        \end{scope}
%
%================================================
%                        
        % polar coordinate system
        \begin{scope}[xshift=7.1cm,yshift=0cm]
\ifdebug
            % grid lines
            \path[draw,grid=0.25cm] (-3.50,-3.00) grid (3.00,2.75);
\fi

            \begin{scope}[rotate around={27.5:(0,0)}]
                
                % objects
                \path coordinate (O) at (0,0);          % origin O
                \path coordinate (P) at (-1.75,{-sin(1.75*180/pi)*1.75*1.1});      % point P

                %  arc from origin to point P
                \draw[scale=1.0, domain=-1.75:0, smooth, samples=1000, gray, line width=1mm, variable=\x]
                    plot ({\x}, {-sin(\x*180/pi)*\x*1.1});
                
                % trajectory
                \draw[scale=1.0, domain=-3.25:0, smooth, variable=\x] plot ({\x}, {-sin(\x*180/pi)*\x*1.1});
                \draw[scale=1.0, domain=0:3.05, smooth, variable=\x] plot ({\x}, {sin(\x*180/pi)*\x*1.1});
                % indicate positive and negative direction along the curve
                \path node at (  3.10,{ sin(3.10*180/pi)*3.10*1.1}) {$\mathbf{+}$};
                \path node at ( -3.30,{-sin(3.30*180/pi)*3.30*1.1}) {$\mathbf{-}$};
                
                % origin
                \path node at (O) {$\mathcal{O}$};
                % tick marks at unit intervals
%                 \foreach \x in {-3,...,2}
                \path[tick] (0,-2pt) -- (0,+2pt);
                % point P
                \filldraw[blue] (P) circle (2pt);
                \path node[text=blue,below right, inner sep=1pt]
                    at (P) {$\mathcal{P}=(d_\mathcal{P})$};
            
            \end{scope}
            
            % label sub figure
            \path node[above right] at (-3.55,-3.00) {[\textbf{b}]};
                    
        \end{scope}
%
%================================================
%                    
    \end{tikzpicture}
    
    \caption[One-dimensional coordinate systems]
            {One-dimensional coordinate systems
            \begin{description}%[nosep,left=\parindent,itemindent=-20pt]
                \item[{[a]}] Straight line \showthe\labelwidth
                \item[{[b]}] Curved path
            \end{description}
            }
            
    \label{F:CS_1D}
    
\end{figure}